% Separated candidate v2; frozen v1 and main are unchanged.
% BEGIN SAFE-PROOF
\begin{proof}
Partition the common fibre graph into its fixed-middle layers $H_v$
and the edge set $H_*$ changing the middle coordinate. In any block,
if $m$ edges remain outside $P_t\cup P_s$ and the two cycles retain
$x,y$ of them, they share at least $\max(0,x+y-m)$ safe edges.
These bounds add over disjoint blocks. Appendix~\ref{app:path232-certificates}
classifies all allowed transitions and gives sufficient block counts
for seven families, each yielding at least two.

Let $a,b,c,d$ be distinct and use the selected cycles defined there.
For $aab\to bab$, the common middle-$b$ prism has no special edges.
The $E$ cycle retains one edge in each row triangle and two vertical
edges. The cycle $M_0$ contains both these triangle edges, and $M_1$
both these vertical edges, giving two safe common edges in either case.

For $aba\to cba$, $H_a$ has nine edges. Both cycles omit
$bab$--$bac$ and $bab$--$bad$: the $M$ layer isolates $bab$, while
the $A$ path starts $bab,dab$. They retain three and five edges in
this seven-edge universe, giving one common edge. In $H_c$ they
retain three and two of four edges, giving another. Both special
sets are empty. Colour permutations and reversal transport all
arguments, proving the assertion for every allowed labelled pair.
\end{proof}
% END SAFE-PROOF

% BEGIN SELECTED-SUPPORT
\subsection{Selected-layer identities and capacities}
Order the end-, middle-, and all-distinct types as $E<M<A$.
For distinct $a,b,c,d$, let $E(aab),M_0(aba),A(abc)$ be the images of
$D_E,D_M,D_A$ under $(0,1,2,3)\mapsto(a,b,c,d)$, and let $M_1$ be
the reversal of $M_0$. Unused-colour interchange fixes $E,M_0,M_1$;
reversal followed by endpoint-colour interchange fixes $A$. Thus
$E,A$ each have one allowed edge set and $M$ has the two stated sets.
Up to these symmetries, the good transitions to a type no smaller are
\[
E(aab):aac,abb,acb,bab,cab;\quad
M(aba):aca,cba;\quad A(abc):abd,adc.
\]
These exhaust one-coordinate changes; bad targets and smaller types
are excluded, and unused colours and the two ends are identified
by the stated symmetries.

In $E(aab)$, the middle-$b$ layer is $2P_3$ and the middle-$c,d$
squares are $P_4$. In $M_0(aba)$, the middle-$a$ layer is $P_8$
with isolated $bab$, and the middle-$c,d$ squares are $P_4$;
reverse for $M_1$. In $A(abc)$, the middle-$a,c$ layers are $P_6$
and the middle-$d$ square is $P_3$ with isolated $bdb$.
These are subgraphs of the respective selected cycles.

Each cycle has four middle-changing edges. For $E$, middle colour $b$
joins $c,d$ twice each, with last coordinate $a$ twice and $c,d$
once each. For $M_0$, middle colour $a$ joins $c,d$ twice each and
every edge has first coordinate $b$; reverse for $M_1$. For $A$,
middle colour $d$ joins $a,c$ twice each, with first coordinate in
$\{b,c\}$ and last coordinate in $\{a,b\}$. These paths and counts
are read directly from~\eqref{eq:path232-cycles}.

For common endpoint lists $L,R$, the layer $H_v$ is
$K_{|L\setminus\{v\}|}\square K_{|R\setminus\{v\}|}$, and there are
$|L\setminus\{v,w\}|\,|R\setminus\{v,w\}|$ edges between middle
colours $v,w$. Removing the special edges gives the safe universe
size $m$; the original cyclic words give the retained counts $x,y$.
The following seven families supply the capacities used in the proof:
\[
\begin{array}{c|c|c}
t\to s&\text{blocks}&(m,x,y)\\\hline
aab\to aac&H_*&(4,3,3)\\
aab\to abb&H_c&(2,2,2)\\
aab\to acb&H_*&(2,2,2)\\
aab\to cab&H_c,H_d&(3,2,2),(1,1,1)\\
aba\to aca&H_d&(4,3,3)\\
abc\to abd&H_*&(6,4,4)\\
abc\to adc&H_a,H_c&(9,5,5)\text{ each}
\end{array}
\]
% END SELECTED-SUPPORT
